\documentclass[12pt]{article}
\usepackage{amsmath, amssymb}
\usepackage{tikz}
% \usetikzlibrary{arrows.meta, positioning}
\usepackage{graphicx} % Required for inserting images
\usepackage[backend=biber, style=authoryear]{biblatex}
\addbibresource{references.bib}

\title{What does the VLM See?}
\author{ML Team, Phronetic AI}
\date{June 2026}

\begin{document}
\maketitle

\begin{abstract}
    Video Large Language Models process all visual tokens uniformly across every transformer layer, regardless of their contribution to the final answer. The resulting inference cost scales quadratically with sequence length, making deployment under tight compute budgets challenging. Existing token reduction methods operate at one of two levels: frame selection before the model, or attention-based pruning inside the model. Frame-level methods such as VideoITG (\cite{wang2026videoitgmultimodalvideounderstanding}) and PEEK (\cite{steunou2026peekpickingessentialframes}) require either a text query at selection time or reference caption supervision, limiting their generality. Token-level methods such as  (\cite{huang2026kitokekernelbasedintervalawaretoken}, DyCoke (\cite{tao2025dycokedynamiccompressiontokens}, either operate without any language grounding (remaining agnostic to what the model actually needs), or derive supervision from task-specific reward signals that are expensive to compute.
    We take a different position: the frozen VLM itself contains a reliable, model-native supervision signal for token importance, localisable to a small subset of intermediate layers. Recent mechanistic analysis of Video-LLMs (\cite{gou2025empiricalstudyvideollmsanswer}) demonstrates that language-to-video attention at layers 12–16 accounts for the disproportionate majority of visual information retrieval, while temporal and spatial self-attention between video tokens contributes comparatively little. Crucially, this signal is distinct from the CLS or BOS attention exploited by earlier pruning methods (such as FastV), which are known to exhibit positional bias and inflate importance scores for semantically empty tokens (\cite{myrzakhan2026sinkawarepruningdiffusionlanguage}). Language-to-video attention at critical intermediate layers does not exhibit this bias \cite{takezoe2026learnprunerrethinkingattentionbasedtoken}.
    We train a lightweight MLP scorer (~1.0M parameters) online, supervised by language-to-video attention extracted from layers 12–16 of a frozen Qwen2.5-VL during forward passes on standard VideoQA data. The scorer reads the merged visual tokens the language model itself consumes and gates them before the decoder. To prevent temporal coverage collapse - where globally sorting top tokens concentrate in a single visual moment, we apply stratified selection across non-overlapping temporal bins ensuring the retained set spans the full video. Spatial and temporal identity of the token is preserved via RoPE, so the model always knows where and when each token originated from. Only the scorer is trained; VLM remains frozen.

    Our contributions are: (1) a training signal derived from model-internal language-to-video attention at empirically validated critical layers, requiring no caption alignment, no reward model, and no architectural modification; (2) input-level token gating that prunes the merged visual tokens before the language-model decoder, saving full-depth LLM compute rather than pruning at an intermediate layer as attention-based methods do; (3) stratified temporal selection ensuring coverage-aware retention without merging or architectural overhead; and (4) an empirical finding that norm-based heuristics reliably identify uninformative tokens but fail to identify important ones, directly motivating the need for a learned selector. This work extends our prior result (\cite{dave2026adaptivetokenisationtemporalredundancy}) that motion-based heuristics suffice for video compression but are structurally insufficient for understanding
\end{abstract}

\section{Introduction}


The dominant paradigm in video understanding has shifted decisively toward large multimodal models. Video Large Language Models (Video-LLMs) now achieve strong performance across a wide range of tasks such as video question answering, temporal reasoning, dense captioning done by processing sampled frames as sequences of visual tokens alongside a language query. Yet this capability comes at substantial computational cost. The quadratic scaling of attention with sequence length means that visual tokens alone account for the overwhelming majority of inference FLOPs, often exceeding language token compute by two orders of magnitude (\cite{gou2025empiricalstudyvideollmsanswer}).

The standard response to this bottleneck has taken two forms. The first is frame-level selection: choose which frames to show the model before any encoding occurs. Methods in this family, including PEEK (\cite{steunou2026peekpickingessentialframes}) and VideoITG (\cite{wang2026videoitgmultimodalvideounderstanding}), have demonstrated that adaptive frame selection outperforms uniform sampling, particularly at low frame budgets. However, they operate on decoded pixels via external vision-language models, their supervision is tied to reference captions or text queries available only at selection time, and they discard entire frames rather than making fine-grained spatial decisions within frames.

The second response is token-level pruning inside the model: retain all frames but reduce the number of tokens that propagate through the transformer layers. Training-free methods such as KiToke (\cite{huang2026kitokekernelbasedintervalawaretoken}) and DyCoke (\cite{tao2025dycokedynamiccompressiontokens}) reduce spatiotemporal redundancy via diversity maximisation or cosine similarity thresholding. These approaches are principled and require no training, but they are query-agnostic, optimising for visual diversity without reference to what the model actually attends to when answering a question. Training-based methods such as (\cite{wang2026dynamictokencompressionefficient}) learn a selection policy via reinforcement learning, using task accuracy as a reward signal. This produces a task-grounded selector but at substantial training cost: each update requires multiple full VLM forward passes to evaluate sampled token masks.

We argue that both families miss a simpler and more direct source of supervision: the VLM itself. Recent mechanistic analysis \cite{gou2025empiricalstudyvideollmsanswer} shows that language-to-video attention at a small subset of intermediate layers, specifically layers 12–16 in 3B to 7B-scale models which has shown to account for the disproportionate majority of visual information retrieval in VideoQA. Most other layers contribute negligibly; temporal and spatial self-attention between video tokens is largely redundant. This finding implies that the model already "knows" which tokens matter, and that knowledge is readable from its own attention patterns at a well-defined depth.

We exploit this directly. We train a lightweight MLP scorer ($\sim$1.0M parameters) online, supervised by language-to-video attention extracted from layers 12–16 of a frozen Qwen2.5-VL. The scorer operates on the merged visual tokens — those the language model actually consumes — gating them before any LLM decoder computation occurs. Crucially, we address a failure mode common to global top-k selection: without temporal constraints, high-scoring tokens can concentrate in a single visually rich moment, leaving later frames entirely unrepresented. We resolve this with stratified temporal selection, distributing the token budget across non-overlapping temporal bins while the scorer determines which tokens within each bin to retain. Spatial and temporal positional identity is preserved through M-RoPE position IDs, so the VLM retains full knowledge of where and when each surviving token originated.

This work connects directly to our prior contribution (\cite{dave2026adaptivetokenisationtemporalredundancy}), which demonstrated that simple L1 delta thresholding over continuous video latents suffices for high-quality reconstruction without learned routing. That result showed motion-based heuristics are sufficient for compression. Here we show they are insufficient for understanding: norm-based heuristics reliably discard uninformative tokens (bottom-k selection is stable) but fail to identify semantically important ones (top-k selection degrades prediction consistency). This asymmetry directly motivates the learned scorer proposed here.

\section{Method}
 
Our method consists of three components: a lightweight token scorer
trained online on a frozen VLM's internal attention, a stratified
temporal selection strategy that enforces coverage across the video,
and an input-level gating mechanism that removes low-scoring tokens
before any LLM computation occurs.
 
\subsection{Preliminaries}
 
A Video-LLM typically consists of a frozen visual encoder
$\text{Enc}_v$, a projector $\text{Proj}$, and a frozen decoder-only
language model with $L$ transformer layers. Given a video $V$ sampled
at $T$ frames, the visual encoder produces patch embeddings:
%
\begin{equation}
    F = \text{Enc}_v(V) \in \mathbb{R}^{T \times N \times d_v}
\end{equation}
%
where $N$ is the number of patches per frame and $d_v$ is the visual
encoder hidden dimension. These are projected into the LLM token space
via $\text{Proj}$, concatenated with text tokens, and processed
through $L$ transformer layers. The full sequence length $T \cdot N +
T_{\text{text}}$ drives the quadratic attention cost, motivating
reduction of $T \cdot N$ before any LLM computation.
 
\subsection{Supervision Signal: Language-to-Video Attention}
 
We derive per-token importance scores from the frozen VLM's internal
attention at a set of critical layers $\mathcal{L} = \{12, 13, 14,
15, 16\}$, identified empirically by \textcite{gou2025empiricalstudyvideollmsanswer} as the
layers where language-guided visual retrieval is concentrated.
 
At each critical layer $\ell \in \mathcal{L}$, the attention module
produces a weight matrix $A^{(\ell)} \in \mathbb{R}^{H \times S \times
S}$, where $H$ is the number of attention heads and $S = T \cdot N +
T_{\text{text}}$ is the full sequence length. We extract the
language-to-video slice i.e., query positions corresponding to text
tokens, key positions corresponding to video tokens and aggregate:
%
\begin{equation}
    s_i = \frac{1}{|\mathcal{L}|}
    \sum_{\ell \in \mathcal{L}}
    \frac{1}{H \cdot T_{\text{text}}}
    \sum_{h=1}^{H}
    \sum_{j=1}^{T_{\text{text}}}
    A^{(\ell)}_{h,\; T \cdot N + j,\; i}
\end{equation}
%
This yields a scalar importance score $s_i$ for each video token $i
\in \{1, \ldots, T \cdot N\}$. Scores are min-max normalised per
sample to $[0, 1]$. During training, the VLM forward pass is truncated
at layer 16 and attention weights are detached; no gradients flow into
the VLM.
 
We use text-to-vision attention specifically because it is resistant to
the positional bias and attention sink effects that afflict vision
encoder CLS attention and early-layer LLM
attention. Vision-to-vision
attention exhibits consistent recency bias; text-to-vision attention in
middle layers does not (\cite{takezoe2026learnprunerrethinkingattentionbasedtoken}).
 
\subsection{Token Scorer}
 
The scorer is an MLP with a LayerNorm input, two GELU hidden layers
(each followed by dropout $0.1$), and a linear head, mapping each merged
visual token to a scalar importance logit:
%
\begin{equation}
    \hat{s}_i = w_o^{\top}\,\sigma\!\big(W_2\,\sigma(W_1\,\text{LN}(f_i)
    + b_1) + b_2\big) + b_o,
    \qquad \sigma = \text{GELU},
\end{equation}
%
where $f_i \in \mathbb{R}^{d}$ is the merged visual token at position
$i$ --- the post-merger representation the language model consumes, of
width $d = 3584$ for the 7B base ($2048$ for the 3B) --- aligned
one-to-one with the language-to-video teacher scores. Here $W_1 \in
\mathbb{R}^{d_h \times d}$, $W_2 \in \mathbb{R}^{d_h \times d_h}$, $w_o
\in \mathbb{R}^{d_h}$, and $d_h = 256$, giving approximately $1.0$M
trainable parameters for the 7B base ($0.6$M for the 3B). The frozen
visual encoder, merger/projector, and all LLM layers are never updated.
 
\paragraph{Training Objective.}
Token selection is fundamentally a ranking problem: only the relative
order of importance scores determines which tokens survive, not their
absolute values. We therefore adopt the ListMLE list-wise ranking
loss (~\cite{10.5555/3020751.3020798}). Let $\pi^*$ denote the permutation of
token indices sorted by decreasing teacher score $s_i$. The loss is:
%
\begin{equation}
    \mathcal{L}_{\text{ListMLE}} =
    -\sum_{i=1}^{T \cdot N}
    \log \frac{
        \exp(\hat{s}_{\pi^*(i)})
    }{
        \displaystyle\sum_{\tau=i}^{T \cdot N}
        \exp(\hat{s}_{\pi^*(\tau)})
    }
\end{equation}
%
This is the Plackett-Luce negative log-likelihood of observing the
teacher-induced ranking under the predicted logits, optimising the
full ranking jointly rather than independently per token.
 
\paragraph{Online Training.}
At each training step, a video-question-answer triple is sampled from a
VideoQA dataset. A single forward pass through the frozen VLM to layer
16 with \texttt{output\_attentions=True} produces the supervision
signal $s$. The scorer forward pass produces $\hat{s}$. The ListMLE
loss is computed and gradients are propagated through the scorer only.
No labels beyond the standard VideoQA annotations present in the
dataset are required.
 
\subsection{Stratified Temporal Selection with Pareto-Adaptive Budgeting}

A global top-$k$ selection strategy applied to scorer logits risks temporal coverage collapse: tokens from visually rich frames accumulate disproportionately, leaving later or quieter frames entirely unrepresented. Conversely, a uniform per-bin allocation $k_t = \lfloor K / T \rfloor$ wastes token budgets on static, uninformative background frames. 

We resolve this by setting $k_t$ adaptively, exploiting the empirical observation that the tail of visual token importance follows a Pareto distribution, whose shape parameter $\alpha$ dictates the concentration of information. To properly estimate this, we must first establish the optimal estimator for a Pareto tail.

\paragraph{The Naïve Model and its Maximum Likelihood Estimator.}
Assume a set of token norms $x_i \overset{\text{iid}}{\sim} \text{Pareto}(\alpha, x_m)$ with density $f(x) = \frac{\alpha\, x_m^\alpha}{x^{\alpha+1}}$ for $x \ge x_m$. The joint likelihood function for $n$ observations is the product of their individual densities:
%
\begin{equation}
    L(\alpha) = \prod_{i=1}^n \frac{\alpha\, x_m^\alpha}{x_i^{\alpha+1}}
\end{equation}
%
Taking the natural logarithm transforms this product into a sum, giving the log-likelihood:
%
\begin{equation}
    \ell(\alpha) = \sum_{i=1}^n \log \left( \frac{\alpha\, x_m^\alpha}{x_i^{\alpha+1}} \right) = n\log(\alpha) + n\alpha\log(x_m) - (\alpha + 1)\sum_{i=1}^n \log(x_i)
\end{equation}
%
To find the maximum likelihood estimator (MLE), we take the derivative with respect to $\alpha$ and set it to zero:
%
\begin{equation}
    \frac{\partial \ell}{\partial \alpha} = \frac{n}{\alpha} + n\log(x_m) - \sum_{i=1}^n \log(x_i) = 0
\end{equation}
%
Solving for $\alpha$ and pulling the $n\log(x_m)$ term into the summation yields the exact MLE:
%
\begin{equation}
    \hat{\alpha}_{\text{MLE}} = \frac{n}{\sum_{i=1}^n \log(x_i / x_m)}
\end{equation}
%
This formulation is statistically robust because the transformation $Y = \log(X/x_m)$ yields an Exponential distribution. By evaluating the Pareto CDF at $x_m e^y$, we find $F_Y(y) = 1 - (x_m / x_m e^y)^\alpha = 1 - e^{-\alpha y}$, which is exactly $Y \sim \text{Exp}(\alpha)$. Consequently, the MLE for $\alpha$ inherits the standard properties of an exponential rate estimator, including a known exact variance and inverse-Gamma sampling distribution.

\paragraph{Extreme Value Theory and the Hill Estimator.}
However, directly applying $\hat{\alpha}_{\text{MLE}}$ to the full distribution of ViT activation norms commits a model misspecification. The bulk of the distribution is approximately log-normal; only the extreme upper region—where the semantically important tokens live—is Pareto. Fitting the naïve model to the whole sample allows the bulk to dominate, resulting in inflated and unstable $\alpha$ estimates.

Following Extreme Value Theory, we relax the assumption to a regularly varying tail: $\Pr(X > x) = x^{-\alpha} L(x)$, where $L(x)$ is a slowly varying function. Given a sufficiently high threshold $u$, the scaled exceedances $X/u$ are approximately $\text{Pareto}(\alpha, 1)$. We can therefore adapt our derived MLE to operate strictly conditionally on the top $k$ order statistics. Substituting $n$ with $k$, the threshold $x_m$ with $x_{(n-k)}$, and evaluating over the top $k$ values $x_{(n-i+1)}$, we obtain the conditional MLE, whose inverse is the well-known Hill estimator:
%
\begin{equation}
    \hat{H}_{k,n} = \frac{1}{\hat{\alpha}} = \frac{1}{k}\sum_{i=1}^{k} \log\frac{x_{(n-i+1)}}{x_{(n-k)}}
\end{equation}

\paragraph{Rényi Representation and Finite-Sample Properties.}
To rigorously analyze the bias-variance trade-off in selecting the threshold $k$, we must establish the exact finite-sample distribution of the Hill estimator. As established previously, for an exact Pareto tail, the transformation $Y = \log(X/x_{(n-k)})$ yields $Y \sim \text{Exp}(\alpha)$. Because $\alpha$ acts as a rate parameter, we can express this as a scaled standard exponential: $Y = E/\alpha$, where $E \sim \text{Exp}(1)$. Substituting this scaling into the Hill estimator translates the components into differences of standard exponential order statistics:
%
\begin{equation}
    Y_{(n-i+1)} - Y_{(n-k)} = \frac{1}{\alpha} \big( E_{(n-i+1)} - E_{(n-k)} \big)
\end{equation}
%
Here, we invoke the Rényi representation theorem. Due to the unique memorylessness property of the exponential distribution, the excesses of the top $k$ order statistics over the threshold $E_{(n-k)}$ are strictly independent of the threshold itself. Furthermore, these $k$ gaps are jointly distributed exactly as the order statistics of a fresh, independent sample of $k$ standard exponentials, denoted $E'_j \sim \text{Exp}(1)$. Since summing a set of order statistics is mathematically identical to summing the unsorted variables, the Hill estimator simplifies remarkably in distribution:
%
\begin{equation}
    \hat{H}_{k,n} \;\overset{d}{=}\; \frac{1}{\alpha k} \sum_{j=1}^{k} E'_j
\end{equation}
%
This elegant reduction frames the inverse Hill estimator simply as $1/\alpha$ scaled by the sample mean of $k$ iid standard exponential variables. Given $\mathbb{E}[E'_j] = 1$ and $\operatorname{Var}(E'_j) = 1$, we can isolate the exact finite-sample expectation and variance under a strict Pareto tail:
%
\begin{equation}
    \mathbb{E}[\hat{H}_{k,n}] = \frac{1}{\alpha k} \sum_{j=1}^{k} \mathbb{E}[E'_j] = \frac{1}{\alpha}
\end{equation}
%
\begin{equation}
    \operatorname{Var}(\hat{H}_{k,n}) = \frac{1}{\alpha^2 k^2} \sum_{j=1}^{k} \operatorname{Var}(E'_j) = \frac{1}{\alpha^2 k}
\end{equation}
%
This exact variance cleanly highlights the central tension in threshold selection. To drive the variance of the estimator to zero, we require $k \to \infty$. However, as $k$ increases, the threshold $x_{(n-k)}$ descends into the non-Pareto bulk of the distribution, violating the slowly varying assumption of $L(x)$ and introducing massive systematic bias. Thus, temporal budgeting requires finding the optimal plateau where $k$ is sufficiently large to stabilize the $1/(\alpha^2 k)$ variance, yet strictly constrained such that $k/n \to 0$ to avoid bulk contamination.
 
\subsection{Input-Level Gating and Positional Identity}
 
The retained merged tokens $F_{\mathcal{I}} = \{f_i : i \in \mathcal{I}
\} \in \mathbb{R}^{K \times d}$ are passed directly to the frozen
language model. All other tokens are discarded before any decoder layer
runs, saving full-depth LLM compute proportional to $1 - K /
(T \cdot N)$.
 
Critically, each retained token carries its original spatiotemporal
position identifiers namely the frame index $t$, spatial row $x$, and spatial
column $y$, through the LLM via
M-RoPE (~\cite{bai2025qwen25vltechnicalreport}). The VLM therefore retains full
knowledge of where and when each surviving token originated, without
any positional re-indexing. The full inference pipeline is:
%
\begin{equation}
    V
    \xrightarrow{\text{Enc}_v} F
    \xrightarrow{\text{Merger}} \tilde{F}
    \xrightarrow{\text{Scorer}} \hat{s}
    \xrightarrow{\text{Stratified top-}k} F_{\mathcal{I}}
    \xrightarrow{\text{LLM}} \text{Answer}
\end{equation}

\subsection{A Privileged Gradient Oracle: A Diagnostic Study}

Alongside the language-to-video attention teacher introduced above, we
investigated a stronger, label-dependent supervision signal --- a
Grad-CAM-style gradient oracle --- and used it to probe three questions
that any token-selection method must answer: (i) does a small
answer-bearing token subset exist at all; (ii) is token importance
concentrated enough to prune aggressively; and (iii) can that importance
be predicted from the per-token features a scorer actually sees.

\paragraph{Gradient-oracle supervision.}
Let $v_i \in \mathbb{R}^{d}$ denote the merged visual token occupying the
$i$-th video position of the LLM input sequence (the same activation the
scorer ingests), and let $y^\star$ be the gold answer token. With every
VLM parameter frozen, a single backward pass of the answer
log-probability onto the visual activations yields a per-token saliency,
%
\begin{equation}
    g_i = \frac{\partial \log p(y^\star \mid V, Q)}{\partial v_i},
    \qquad
    s^{\text{o}}_i = \mathrm{ReLU}\!\big(\langle g_i,\, v_i \rangle\big),
\end{equation}
%
min--max normalised to $[0, 1]$. Only the visual activations carry
gradient; no model weight is updated. Because $s^{\text{o}}$
differentiates the \emph{true} answer, it is a privileged teacher --- an
upper bound on what any label-free scorer could recover. We train the
scorer against it with a per-token binary cross-entropy,
%
\begin{equation}
    \mathcal{L}_{\text{BCE}} = -\frac{1}{T N}\sum_{i=1}^{T N}
    \Big[ s^{\text{o}}_i \log \sigma(\hat{s}_i)
    + (1 - s^{\text{o}}_i)\log\big(1 - \sigma(\hat{s}_i)\big) \Big].
\end{equation}

\paragraph{Training procedure.}
Only the scorer is trained; the visual encoder, merger, projector, and
language model remain frozen throughout. At each step we draw a single
video--question--answer triple and run one forward pass through the frozen
VLM to obtain the merged visual tokens $\{v_i\}$ and the answer
log-probability, followed by one backward pass onto $\{v_i\}$ to form the
oracle target $s^{\text{o}}$. The same two-layer MLP of
Section~2.3 maps each merged token to a logit $\hat{s}_i = \text{Scorer}(v_i)$,
and is updated by $\mathcal{L}_{\text{BCE}}$ with AdamW under a
cosine-annealed learning rate and gradient clipping; gradients are
accumulated over several samples to emulate a larger effective batch. No
supervision beyond the standard VideoQA answer label is used. Inference is
identical to the main pipeline --- the predicted scores feed the same
stratified, Pareto-budgeted top-$k$ selection of the preceding
subsections. This route thus differs from our primary method in exactly
two places: the teacher signal ($s^{\text{o}}$ in place of
language-to-video attention) and the loss (BCE in place of ListMLE).

\paragraph{An answer-bearing subset exists (ceiling).}
We first ask whether \emph{any} small token subset suffices for the
answer. Retaining the top-$K$ tokens by oracle score (stratified per
frame) and running the multiple-choice forward on the physically
shortened sequence, the oracle attains $0.780$ accuracy at a retention of
$\rho = 0.1$, versus $0.674$ for matched-retention uniform sampling --- a
gap of $+0.106$ over $4{,}996$ NExT-QA videos. A privileged selector
therefore recovers most of the answer from a tenth of the tokens and
clears uniform by a wide margin, confirming that aggressive token
reduction is achievable in principle. Since the oracle keeps the tokens
it was differentiated to favour, $+0.106$ is an upper bound rather than
an achievable score.

\paragraph{Importance is heavy-tailed (concentration).}
We next characterise the \emph{shape} of the oracle distribution, which
governs how aggressively the budget above may be set. Estimating the tail
index with the moment (Einmahl--de Haan) estimator gives a median
$\hat{\xi} = 0.43$ (IQR $[0.31, 0.55]$), with a median $50.4\%$ of tokens
at exactly zero saliency (post-ReLU) and a median per-sample skew of
$7.6$. Importance is thus sharply Pareto-like: roughly half of the merged
tokens are inert and a small heavy tail carries the signal. This
empirically substantiates the Pareto-tail premise underlying our adaptive
budgeting, and makes a clean, balanced binary ``used vs.\ unused'' target
available.

\paragraph{The oracle is not feature-predictable (negative transfer).}
The two results above are encouraging, but a feature-based scorer can
exploit them only if the oracle ranking is a function of the per-token
features it ingests. We tested this directly, correlating the oracle's
per-token ranking against three label-free signals computed from the same
merged tokens: token norm $\lVert v_i \rVert$, query cosine $\max_t
\cos(v_i, q_t)$ (semantic relevance to the question text), and temporal
novelty $1 - \cos(v_i, v_i^{t-1})$. Across $4{,}996$ videos and two
model/frame configurations, all three are statistically indistinguishable
from chance (Table~\ref{tab:oracle-corr}): every Spearman $\rho$ satisfies
$|\rho| \le 0.07$, and the top-quartile recall hovers at the $0.25$ chance
baseline. Query cosine is the only signal above chance, and only
marginally so, while token norm sits \emph{below} chance. The
answer-criticality the oracle captures is therefore largely invisible to
static per-token features; an MLP regressing these features onto
$s^{\text{o}}$ has little to extract, which is exactly the chance-level
recall we observe when training with $\mathcal{L}_{\text{BCE}}$.

\begin{table}[h]
\centering
\begin{tabular}{l r r r r}
\hline
 & \multicolumn{2}{c}{3B, 8 frames} & \multicolumn{2}{c}{7B, 16 frames} \\
Signal & $\rho$ & R@25 & $\rho$ & R@25 \\
\hline
Token norm $\lVert v_i \rVert$              & $-0.027$ & $0.226$ & $-0.067$ & $0.194$ \\
Query cosine $\max_t \cos(v_i, q_t)$        & $+0.033$ & $0.286$ & $+0.036$ & $0.284$ \\
Novelty $1 - \cos(v_i, v_i^{t-1})$          & $-0.013$ & $0.244$ & $-0.010$ & $0.246$ \\
\hline
\end{tabular}
\caption{Oracle vs.\ feature-only signals: median per-video Spearman
$\rho$ between the gradient oracle and each label-free signal, with
top-quartile recall (R@25; chance $= 0.25$), over $4{,}996$ NExT-QA
videos. Every signal is at or near chance, indicating the oracle is not
feature-predictable.}
\label{tab:oracle-corr}
\end{table}

\paragraph{Implication.}
Taken together, the diagnostic delivers a sharp design constraint. A
concentrated, answer-bearing token subset demonstrably exists (ceiling
$+0.106$) and is heavy-tailed (Pareto $\hat{\xi} = 0.43$), so the
stratified selection and Pareto-adaptive budgeting above are well-founded.
But that subset cannot be reconstructed from static per-token features
alone: importance is a property of a token's \emph{cross-modal
interaction} with the query, not of the token in isolation. This is
precisely why we supervise the scorer from the model's internal
language-to-video attention rather than from a label-dependent gradient
oracle, and it cautions that the scorer input must be expressive enough to
reflect that interaction rather than scoring each token independently.

\section{Proposed Experiments}
 
\subsection{Setup}
 
We evaluate on three VideoQA benchmarks spanning different temporal
scales and task types:
 
\begin{itemize}
    \item \textbf{VideoMME}~\parencite{fu2025videommefirstevercomprehensiveevaluation} - 900 videos
    across six domains with short, medium, and long duration splits.
    Primary benchmark for overall VideoQA accuracy.
 
    \item \textbf{MVBench}~\parencite{li2024mvbenchcomprehensivemultimodalvideo} - 4,000 QA
    pairs across 20 temporally-sensitive task types including action
    sequences, object permanence, and motion analysis.
\end{itemize}
 
The base VLM throughout is Qwen2.5-VL-7B-Instruct. The scorer is
trained on the NExTVideo training split~\parencite{nextqa} for 10,000
steps with AdamW ($\text{lr} = 10^{-4}$, weight decay $10^{-2}$),
cosine annealing schedule, and gradient clipping at 1.0.
 
\subsection{Baselines}
 
We compare at token retention ratios $\rho \in \{25\%, 50\%, 75\%\}$
of the full token count $T \cdot N$:
 
\begin{itemize}
    \item \textbf{Full model} - no token reduction, upper bound.
    \item \textbf{Uniform sampling} - spatially and temporally
    uniform subsampling, strong baseline~\parencite{steunou2026peekpickingessentialframes}.
    \item \textbf{L1 delta thresholding} - our prior BMVC heuristic
    adapted to token-level selection, representing the motion-based
    approach.
    \item \textbf{KiToke}~\parencite{huang2026kitokekernelbasedintervalawaretoken} - best
    published training-free token compression at matched retention
    ratios.
    \item \textbf{SCORE}~\parencite{wang2026dynamictokencompressionefficient} - best published
    training-based token compression, representing RL-supervised
    selection.
    \item \textbf{Ours} - learned scorer with stratified temporal
    selection.
\end{itemize}
 
\subsection{Primary Metrics}
 
VideoQA accuracy (multiple-choice top-1) is the primary metric across
all benchmarks. We additionally report wallclock inference speedup
relative to the full model at each retention ratio, measured on a
single Ampere architecture GPU with batch size 1.
 
\subsection{Ablations}
 
\paragraph{Supervision source.}
We ablate the supervision signal by replacing language-to-video
attention with: (a) uniform random targets, (b) CLS attention from the
vision encoder, and (c) all-token average attention. This directly
tests whether the choice of critical-layer language-to-video attention
is the source of any gains over baselines.
 
\paragraph{Layer range.}
We vary $\mathcal{L}$ across $\{1\text{--}5\}$, $\{6\text{--}11\}$,
$\{12\text{--}16\}$, and $\{17\text{--}28\}$ to verify that the
critical layer range identified by \textcite{gou2025empiricalstudyvideollmsanswer} is
indeed the most informative for our supervision signal.
 
\paragraph{Selection strategy.}
We compare global top-$k$ against stratified top-$k$ to quantify the
benefit of temporal coverage enforcement, particularly on long-form
benchmarks where temporal collapse is most harmful.
 
\paragraph{Norm heuristic vs.\ learned scorer.}
We include a norm-based selector (top-$k$ by L2 norm of ViT features)
as an additional ablation, connecting directly to our empirical finding
that bottom-$k$ norm selection is reliable but top-$k$ norm selection
is not. This motivates the learned scorer within the paper's own
experimental narrative.
 
\subsection{Expected Outcomes}
 
We expect the learned scorer with stratified selection to outperform
all training-free baselines (Uniform, L1 delta, KiToke) at low
retention ratios ($\rho = 25\%$), where the quality of token selection
matters most. The gap over SCORE should emerge from training
efficiency: our method requires a single truncated VLM forward pass per
training step, whereas SCORE requires $K$ full forward passes. At
$\rho = 75\%$, differences between methods are expected to narrow, as
sufficient tokens remain for uniform coverage to be competitive --- a
pattern consistent with findings in \textcite{steunou2026peek} and
\textcite{learnpruner2026}.
 
The ablation on selection strategy is expected to show the largest
gains for stratified over global top-$k$ on EgoSchema, given its
three-minute video length and the higher probability of temporal
collapse at aggressive retention ratios.

\printbibliography
\end{document}
\bibliography{references}
